\documentclass[12pt]{article}
\usepackage[margin=1in]{geometry}
\usepackage{enumitem}
\usepackage{titlesec}
\usepackage{parskip}
\usepackage{graphicx}
\usepackage{xcolor}
\usepackage{hyperref}
\usepackage{fontenc}

\title{\textbf{Schmitter 1974 RG Notes}\\
\large Schmitter, Philippe C.. ``Still the Century of Corporatism?'' [1974].\\
\textit{The Review of Politics} 36(1,): 85-131.}
\author{}
\date{}

\begin{document}

\maketitle

``Corporatism can be defined as a system of interest representation in which the constituent units are organized into a limited number of singular, compulsory, noncompetitive, hierarchically ordered and functionally differentiated categories, recognized or licensed (if not created) by the state and granted a deliberate representational monopoly within their respective categories in exchange for observing certain controls on their selection of leaders and articulation of demands and supports.''

Connect to Stokes et al for clientelism of course, but also a broader conversation of state capture / regulatory capture

Can be tied to Dahl's discussion of polyarchy -- this is particularly noncompetitive, even if it (somehow) was inclusive

Corporatism as a distinct alternative to pluralism wrt interest politics -- alternative model. In particular, corporatism involves cooptation of the state's interest group representation by the interests groups themselves

\end{document}
